\section*{Capítulo \ref{ch:b2:chromatography} --- Cromatografia}

\begin{solution}{exo:b2:chromatography:1}
$k = (5.0 - 1.0)/1.0 = 4.0$; fração do tempo na \omterm{def:b2:chromatography:principle}{fase móvel} $1/(1 + k) = 0.20$.
\end{solution}

\begin{solution}{exo:b2:chromatography:2}
$N = 16(6.0/0.30)^2 = 16 \times 400 = 6400$.
\end{solution}

\begin{solution}{exo:b2:chromatography:3}
$H = \qty{250}{mm}/10\,000 = \qty{0.025}{mm} = \qty{25}{\micro m}$.
\end{solution}

\begin{solution}{exo:b2:chromatography:4}
Etanol no sangue: CG (volátil). Proteína: CLAE. Cafeína em uma bebida: CLAE (em água, pouco volátil
sem preparação). Benzeno na gasolina: CG. Açúcares: CLAE (eles se decompõem antes de ferver).
\end{solution}

\begin{solution}{exo:b2:chromatography:5}
$R_s = 2(4.6 - 4.0)/(0.40 + 0.50) = 1.2/0.90 \approx 1.3$: não totalmente separados até a linha de base (abaixo de 1.5).
\end{solution}

\begin{solution}{exo:b2:chromatography:6}
$\sqrt N = 4R_s\,\dfrac{\alpha}{\alpha - 1}\,\dfrac{1 + k_2}{k_2} = 4 \times 1.5 \times 11 \times 1.25 =
82.5$, logo $N \approx 6.8 \times 10^3$.
\end{solution}

\begin{solution}{exo:b2:chromatography:7}
$u_{\mathrm{opt}} = \sqrt{8/2} = \qty{2}{mm/s}$; $H_{\min} = 5 + 2\sqrt{16} = \qty{13}{\micro m}$.
\end{solution}

\begin{solution}{exo:b2:chromatography:8}
O fenol (\ce{OH} polar, ligações de hidrogênio com a água) primeiro, depois o benzeno, depois o tolueno
(um \ce{CH3} a mais, mais apolar). Mais metanol torna a \omterm{def:b2:chromatography:principle}{fase móvel} menos polar: os três eluem mais cedo,
na mesma ordem.
\end{solution}

\begin{solution}{exo:b2:chromatography:9}
$F = (860/400)/(20.0/10.0) = 2.15/2.00 = 1.075$. Amostra: $c_X = (645/410) \times 10.0/1.075 \approx
\qty{14.6}{mg/L}$.
\end{solution}

\begin{solution}{exo:b2:chromatography:10}
$R_s \propto \sqrt N \propto \sqrt L$: $L$ deve ser multiplicado por $(1.5/1.06)^2 \approx 2.0$, o que
dobra o tempo de análise e a pressão. Outras alavancas: a seletividade (\omterm{def:b2:chromatography:principle}{fase móvel}, \omterm{def:b2:chromatography:principle}{fase estacionária},
temperatura), partículas menores, a vazão ótima.
\end{solution}

\begin{solution}{exo:b2:chromatography:11}
Os picos estão separados por $4\sigma R_s$, de modo que o ponto médio fica a $2\sigma R_s$ de cada um. $R_s = 1$: cada pico
contribui com $\mathrm e^{-(2)^2/2} = \mathrm e^{-2} = 0.135$; o vale fica a $0.27$ da altura do pico.
$R_s = 1.5$: $2\mathrm e^{-(3)^2/2} = 2\mathrm e^{-4.5} \approx 0.022$, cerca de 2~\%: de volta à linha de
base para fins práticos.
\end{solution}

\begin{solution}{exo:b2:chromatography:12}
$k_2/(1+k_2)$: 0.67, 0.83, 0.91; tempo em unidades de $t_M$: 3, 6, 11. Passar de 2 a 5 ganha 25~\% em
\omterm{def:b2:chromatography:separation}{resolução} pelo dobro do tempo; de 5 a 10, 9~\% por quase o dobro outra vez. Um $k$ de cerca de 2 a 5 é o
compromisso habitual.
\end{solution}

\begin{solution}{pb:b2:chromatography:1}
\textbf{1.} Estacionária: grãos de sílica que levam longas cadeias alquila, apolar. Móvel: água com
metanol, polar.
\textbf{2.} Eles são muito polares e ficam na \omterm{def:b2:chromatography:principle}{fase móvel}: $k \approx 0$.
\textbf{3.} A cafeína tem um grupo metila a mais que a teofilina, e a teofilina um \ce{N-H} que forma
ligações de hidrogênio com a água: a teofilina é mais polar e elui primeiro.
\textbf{4.} Ela tem estrutura próxima (resposta e comportamento semelhantes), está ausente da bebida, e
elui perto da cafeína sendo resolvida dela.
\textbf{5.} Ambos diminuiriam: uma \omterm{def:b2:chromatography:principle}{fase móvel} menos polar compete melhor com a \omterm{def:b2:chromatography:principle}{fase estacionária}.
\textbf{6.} O tempo que a \omterm{def:b2:chromatography:principle}{fase móvel} leva para atravessar a coluna, o tempo que um composto passa em movimento.
\textbf{7.} $k_{\mathrm{teo}} = (3.0 - 1.0)/1.0 = 2.0$; $k_{\mathrm{caf}} = 2.4$.
\textbf{8.} $\alpha = 2.4/2.0 = 1.2$.
\textbf{9.} Cafeína $N = 16(3.4/0.20)^2 = 4624 \approx 4.6 \times 10^3$; teofilina $16(3.0/0.18)^2
\approx 4.4 \times 10^3$.
\textbf{10.} $H = \qty{150}{mm}/4624 \approx \qty{0.032}{mm} = \qty{32}{\micro m}$.
\textbf{11.} $R_s = 2(3.4 - 3.0)/(0.18 + 0.20) = 0.80/0.38 \approx 2.1$.
\textbf{12.} $\frac{\sqrt{4624}}{4} \times \frac{0.2}{1.2} \times \frac{2.4}{3.4} = 17 \times 0.167 \times
0.706 \approx 2.0$; a pequena diferença vem da hipótese de larguras iguais da equação.
\textbf{13.} Sim: $R_s > 1.5$.
\textbf{14.} $N$ é dividido por dois: $R_s \approx 2.1/\sqrt2 \approx 1.5$, justo na separação até a linha de base.
\textbf{15.} Dividido por dois: a cafeína em \qty{1.7}{min}.
\textbf{16.} $R_s \approx 2.1\sqrt2 \approx 3.0$, pelo dobro do tempo e o dobro da pressão: um desperdício aqui.
\textbf{17.} A composição da \omterm{def:b2:chromatography:principle}{fase móvel} (fração de metanol, pH, outro solvente orgânico) ou a
\omterm{def:b2:chromatography:principle}{fase estacionária}.
\textbf{18.} Pouco: $k_2/(1 + k_2) = 0.71$ já; com $k_2 = 5$ seria 0.83 (+18~\%) por um
tempo de análise multiplicado por $6/3.4 \approx 1.8$.
\textbf{19.} $F = (1500/1000)/(50.0/40.0) = 1.50/1.25 = 1.20$.
\textbf{20.} Os dois picos vêm da mesma injeção: o volume se cancela na razão das áreas.
\textbf{21.} $c = (970/1010) \times 40.0/1.20 \approx \qty{32.0}{mg/L}$.
\textbf{22.} Diluição de dez vezes: \qty{320}{mg/L}.
\textbf{23.} $A_{IS}$ seria grande demais, e o resultado para a cafeína baixo demais.
\textbf{24.} Uma série de padrões de cafeína injetados exatamente nas mesmas condições que a amostra, uma
reta de calibração da área em função da concentração, e volumes de injeção reprodutíveis.
\textbf{25.} $\qty{320}{mg/L} \times \qty{0.250}{L}$: $\boldsymbol{\approx \qty{80}{mg}}$ de cafeína por lata (dados do
exercício).
\end{solution}
